\chapter{Compatibility of the Mathematical Constructions}
\label{ch:ch13_compatibility}

The question of this chapter is how the results of \cref{ch:ch01_clifford}--\cref{ch:ch12_causal} fit together. The earlier chapters
establish, one sector at a time, what the algebra fixes, what it admits, what it obstructs and what remains open. Here the sectors are held
against one another. \Cref{sec:assessment} sorts the results into the four classes of \cref{sec:deflation} and lists the choices that the algebra
leaves free. \Cref{sec:chains} follows five research chains, from Clifford data to causal evolution, to the physical photon, to interactions,
to stable states and to gravitational dynamics, and marks every link as established, constructed, conditional, obstructed or open.
\Cref{sec:sectors} records for each sector its module, Lorentz action, complex structure, projector, bilinear form, action and status, and
keeps the projectors apart. \Cref{sec:incompatible} names the readings that exclude one another and the bridges that are missing. The last two
sections assess, once, the hypothesis on the rest energy that motivated part of the programme (\cref{sec:resynthesis}) and its second question,
duplication and the $b$-field between manifolds (\cref{sec:dupsynthesis}).

The chains are research chains, not completed derivations. No common action for the sectors is derived or claimed. Where two sectors are
realised together, the joint realisation is stated with its scope; everywhere else they are separate realisations on the same algebra.

% --- source: tex/s31_assessment#0 ((head))


% --- source: tex/s31_assessment#1 (What the algebra decides)
\section{What the algebra decides}
\label{sec:assessment}

We sort the results of this manuscript into the four classes announced in \cref{sec:deflation}, and then name the choices that remain. The test applied throughout is the
deflationary one of \cref{cor:nopredict}: a statement about the arena carries information only if it could have come out otherwise, and compatible representation
data is not a prediction.

% --- source: tex/s31_assessment#2 (Fixed by the algebra)
\subsection{Fixed by the algebra}

At one constituent the algebra fixes more than one might expect. The signature rule separates the bivectors that generate a compact $U(1)$, the gauge-admissible ones, from those that
generate $\R$: as operators on $S$ the first are complex structures and the second have real $\pm1$ eigenspaces (conjugation by either kind is an
involution of the algebra). The
pole operators grade the algebra by $\Z_2^3$ through a maximally split Cartan subalgebra. The pairing of the two triples has exactly two square roots, the grading root
$\cK$ and the binding root $G$, related by the Cartan involution $\Ad_{e_1e_2e_3}$. Under permutation symmetry the clock is the barycentre of the Fisher--Rao octant, a description that is chosen; all unit clocks are related
by rotations, and conversion along any single clock is invisible in the equations of motion, so no clock is selected dynamically. The imaginary unit is selected, not adjoined: inside the joint commutant of the interior
Lorentz action and the internal bivectors the only complex structures commuting with the gamma matrices are $\pm J$. Charge conjugation is the surplus plane, time reversal
squares to $-\id$, the kinetic forms alone fix the internal symmetry algebra $\uu1_\Omega\oplus\su2$ of the free massless action, which becomes a gauge algebra only once a
connection is introduced, the mass form cuts it to $\uu1_J$, and that $\uu1$ is the whole internal
centraliser. Along the tower the internal centraliser is $\so{3n-1,3n-3}$, nested and over one Lorentz frame, and the diagonal Lorentz action gives a compact internal algebra
forming a reductive dual pair with the Lorentz algebra, with the exact identity $C_{\rm cop}+C_\Delta=\tfrac12n(n+2)$; both the dual pair and the identity
are inherited \cite{Howe89,Neergard19}.

In the field-strength reading the centraliser of $G$ is a Lorentz algebra under which $\R^{3,3}$ \emph{is} $\Lambda^2\R^{1,3}$, $G$ is the Hodge star, the two Maxwell
invariants are the real and imaginary parts of one holomorphic quadratic form, and covariance fixes the form of Maxwell's equations through a multiplicity count. In the
interior dynamics, factorisation and binding with $G^2=-\id$ coexist exactly in neutral signature, and, within the stated equivariant one-generator ansatz, the force law closes in one way only, into a rotor acting by $G$ whose free motion is a
pendulum in the duality angle. In the curved arena the curvature operator lives on the frame, the first Bianchi identity is the single trace $\Tr(G\cR)=0$, and the duality
rotation of the curvature ends, where conservation and variationality both select it, at the double dual contracting to the Einstein tensor; in first-order form the same
rotation is a Barbero--Immirzi angle. Kinematically, choosing a Lorentzian four-plane is choosing its definite orthogonal two-plane, represented by a normalised simple bivector
whose two signs encode an orientation, which is a complex structure, and at one constituent the light cone is the set of
chiral spinors squared with the help of $J$.

% --- source: tex/s31_assessment#3 (Constructed)
\subsection{Constructed}

The constructions are existence results; in each case the algebra admits the object and nothing selects it among alternatives. At three constituents the axial projector
$P_*$ of rank $64$ carries the Pati--Salam algebra and one chiral family with a right-handed neutrino, and settles the question of colour with chirality. The diagonal Dirac
block carries a gauge-invariant real $\Cl(1,3)$ built from a copy-singlet cubic. At the mixed second rung the univectorial potential, a Dirac block and the frame field can
be coupled into two-flavour spinor electrodynamics. On the $272$ mass kernels of the chiral block, with the complex structure $\widetilde{\mathcal C}$, the renormalizable potential of the $(\mathbf{10},\mathbf1,\mathbf3)$
component is classified, and sufficient conditions and separate necessary conditions under which the Standard-Model orbit is the global
minimum of that restricted tensor potential are found (\cref{thm:select0,thm:selectk}). In the full $272$-dimensional leading-logarithm model
(an $O(272)$-symmetric tree potential at the matching scale, one gauge and one Yukawa coupling) the radiative spectrum of all $271$ angular modes is
computed and the Standard-Model orbit is a local minimum below $x=27$ (\cref{thm:fullsecond}); its global minimality below $x_*$ is a conjecture supported
by numerics (\cref{prop:randomstart}), and \cref{thm:compat} gives compatibility and spectra at the chosen vacuum, not full-space selection. For a given bidoublet vacuum value the electroweak
stage of \cref{prop:EW}, the seesaw and the Georgi--Jarlskog factor follow. At the second rung a Lorentz-invariant, uncharged, non-rigid
vacuum with a charged-scalar current exists on a compact $S^4$ of projectors, with a curved extension in which it is parallel. Both readings of $\R^{3,3}$ construct the
Einstein tensor of a given four-metric.

% --- source: tex/s31_assessment#4 (Obstructed)
\subsection{Obstructed}

Each obstruction holds under hypotheses stated where it was proved, and those hypotheses are part of the result. Hypercharge cannot be read off anomaly cancellation for
charges built from three commuting involutions; $\su3$ fits in the effective algebra of no axial projector of the five-projector family; the two central $\uu1$s of $P_*$
cannot be gauged on its matter; the diagonal sector is vector-like on its complete blocks. The two readings of the same six directions exclude one another, and the
obvious repair at the second rung costs the internal algebra, whose centraliser shrinks to the rotation of one surplus pair. The positive-form interior oscillates only at zero momentum. No Lie--Poisson closure lives on the frame
alone, and no inertia is both Lorentz invariant and rate-conserving. Non-orthogonality of $\bE$ and $\bB$ breaks $P$ and $T$, is never spontaneous for the Maxwell energy, and generates no mass.
The paired-mode projector sources only a massive vector, and at the first rung a vacuum that is both massless and Lorentz invariant is rigid. The normal $U(1)$ of the curved
model fails as electromagnetism on three independent counts: a spinless particle in the zero-normal-momentum reduced sector cannot see it, its adapted generator does not commute with the physical generators, and it
gravitates with the wrong sign relative to the ambient Einstein--Hilbert action. At the third rung the adapted-generator count is removed by a commuting vacuum, the sign count concerns a metric origin only, and a
fourth count decides: on the chiral block its charge is half the fermion number, with a mixed anomaly with $\su2_L$ (\cref{rem:photonn3}). The electric charge is the
photon candidate there, while the vacuum selected by the adopted potential is electroweak symmetric and the electroweak breaking mechanism is missing. The Ricci--DeTurck method does not transfer to split signature (other split-signature flows are not excluded). No vacuum value of a field built from the internal bivectors and $J$ makes the matter
chiral, Hermitian order parameters keep rank five, and vacuum alignment does not see chirality. In the causal analysis, the six-dimensional free field is ill-posed and has
no time orientation, no $SO(3,3)$-invariant polynomial potential of an unconstrained adjoint scalar of degree below eight isolates the Lorentzian type, every such
invariant potential is degenerate with a Lorentz-violating orbit and has six non-Goldstone flat directions, every factorised invariant kinetic term has ghosts, and,
relative to the ambient Einstein--Hilbert action, every vector from the normal plane is a ghost.

% --- source: tex/s31_assessment#5 (Open, and the choices that remain)
\subsection{Open, and the choices that remain}
\label{sec:openchoices}

\Cref{tab:choices} lists the choices that the algebra leaves free, with the place where each enters. Most of the open problems are downstream of them.

\begin{table}[t]
\centering
\small
\begin{tabular}{@{}>{\raggedright\arraybackslash}p{4.0cm}>{\raggedright\arraybackslash}p{10.4cm}@{}}
\toprule
choice & status\\
\midrule
the projector & a field with a computable configuration space; at rank $64$ at least two admissible orbits, of dimensions $34$ and $20$; a spectrum-only potential is constant on
them, so selection needs couplings that distinguish orbits\\
the Lorentz action & spectator and diagonal actions are not conjugate and trade against each other: a large internal algebra with ghost-signed directions and a chiral block,
or a small compact algebra on vector-like matter\\
the hypercharge direction & a compatible embedding exists on $P_*$ and is output by the vacuum of the kernel field, but the vacuum itself is selected only on an open set of
couplings\\
the binding operator and its reading & $G$ is adopted; conditional uniqueness once its Lorentz algebra is taken, not selection; the field-strength reading and the interior reading
of the same six directions exclude one another\\
scales and inertias & no ratio among the rotor inertia, the frame inertia and the coupling scales is fixed; the inertia of the closure has $21$ free directions satisfying neither
requirement\\
the action & no parent action is derived; the gauge field, the metric and the kernel field are given dynamics by adopted actions whose equations follow\\
the four-plane and positivity & the selector is $J$; the type is not selected by the dynamics of an unconstrained invariant adjoint scalar, and every healthy
construction used a positive inner product\\
families & the block carries one family; three copies of a module are a multiplicity index, not three constituents, and nothing reduces the copy multiplicity to three\\
\bottomrule
\end{tabular}
\caption{The choices the algebra leaves open.}
\label{tab:choices}
\end{table}

The open problems that these choices leave, with the steps that could settle them, are collected in \cref{sec:open}.

\section{Five research chains}
\label{sec:chains}

Each chain runs from algebraic data to a physical conclusion. A link is \emph{established} when it is a theorem of the manuscript or an
inherited result used within its hypotheses; \emph{constructed} when the object exists and is exhibited but nothing selects it;
\emph{conditional} when it holds under an assumption, an adopted action or a chosen datum, stated where it enters; \emph{obstructed} when it is
impossible under stated hypotheses; and \emph{open} otherwise. Where a link has two statuses, for instance existence established and
selection conditional, \cref{fig:chains} draws the weaker one and \cref{tab:chains} gives both. \Cref{fig:chainsdetail} draws the same
chains in detail: what is proved at each stage, with the place where it is proved, the choice or assumption on which each link rests, the
no-go theorems that close routes along each chain, what remains open, and the results that join the chains.

\begin{figure}[t]
\centering
\begin{tikzpicture}[font=\footnotesize,>=Stealth,
  box/.style={draw,rounded corners=2pt,align=center,text width=2.3cm,minimum height=0.95cm,inner sep=2pt},
  est/.style={->,very thick,gcol},
  con/.style={->,very thick,tcol},
  cnd/.style={->,very thick,ocol,dashed},
  obs/.style={->,very thick,scol,densely dotted},
  opn/.style={->,very thick,mgray,dotted},
  lab/.style={font=\scriptsize,above,inner sep=1pt}]
\foreach \r/\name in {0/spacetime,1/electromagnetism,2/matter and interactions,3/energy and mass,4/gravity}
  \node[anchor=west,font=\footnotesize\itshape] at (-1.45,{-1.75*\r+0.78}) {\name};
% row 0: spacetime
\node[box] (a1) at (0,0) {Clifford data};
\node[box] (a2) at (3.85,0) {Lorentz structure};
\node[box] (a3) at (7.7,0) {observable four-plane};
\node[box] (a4) at (11.55,0) {causal evolution};
\draw[cnd] (a1) -- node[lab]{cond.} (a2); \draw[cnd] (a2) -- node[lab]{cond.} (a3); \draw[opn] (a3) -- node[lab]{open} (a4);
% row 1: electromagnetism
\node[box] (b1) at (0,-1.75) {field-strength reading};
\node[box] (b2) at (3.85,-1.75) {charge and gauge geometry};
\node[box] (b3) at (7.7,-1.75) {Maxwell dynamics};
\node[box] (b4) at (11.55,-1.75) {physical photon};
\draw[cnd] (b1) -- node[lab]{cond.} (b2); \draw[cnd] (b2) -- node[lab]{cond.} (b3); \draw[opn] (b3) -- node[lab]{open} (b4);
% row 2: matter
\node[box] (c1) at (0,-3.5) {tower and module};
\node[box] (c2) at (3.85,-3.5) {chiral content};
\node[box] (c3) at (7.7,-3.5) {anomaly-free gauge action};
\node[box] (c4) at (11.55,-3.5) {interactions};
\draw[con] (c1) -- node[lab]{constr.} (c2); \draw[cnd] (c2) -- node[lab]{cond.} (c3); \draw[opn] (c3) -- node[lab]{open} (c4);
% row 3: energy and mass
\node[box] (d1) at (0,-5.25) {geometric operators};
\node[box] (d2) at (3.85,-5.25) {interior and field dynamics};
\node[box] (d3) at (7.7,-5.25) {energy and mass kernels};
\node[box] (d4) at (11.55,-5.25) {stable physical states};
\draw[con] (d1) -- node[lab]{constr.} (d2); \draw[cnd] (d2) -- node[lab]{cond.} (d3); \draw[opn] (d3) -- node[lab]{open} (d4);
% row 4: gravity
\node[box] (e1) at (0,-7.0) {curved geometry};
\node[box] (e2) at (3.85,-7.0) {curvature identity};
\node[box] (e3) at (7.7,-7.0) {variational field equations};
\node[box] (e4) at (11.55,-7.0) {healthy gravitational dynamics};
\draw[est] (e1) -- node[lab]{estab.} (e2); \draw[opn] (e2) -- node[lab]{open} (e3); \draw[obs] (e3) -- node[lab]{obstr.} node[below,font=\scriptsize,inner sep=1pt]{open} (e4);
% legend
\begin{scope}[yshift=-8.3cm,xshift=0.2cm]
\draw[est] (-0.6,0) -- (0.2,0); \node[anchor=west] at (0.25,0) {established};
\draw[con] (2.4,0) -- (3.2,0); \node[anchor=west] at (3.25,0) {constructed};
\draw[cnd] (5.4,0) -- (6.2,0); \node[anchor=west] at (6.25,0) {conditional};
\draw[obs] (8.4,0) -- (9.2,0); \node[anchor=west] at (9.25,0) {obstructed};
\draw[opn] (11.3,0) -- (12.1,0); \node[anchor=west] at (12.15,0) {open};
\end{scope}
\end{tikzpicture}
\caption{The five research chains of \cref{sec:chains}, each link drawn with its weakest status (\cref{tab:chains} gives both statuses where there are two).
The diagram is a summary of the statuses proved or stated in the cited sections; it is not a derivation, and no chain is complete.}
\label{fig:chains}
\end{figure}

\begin{figure}[p]
\centering
\colorlet{ngcol}{scol!55!tcol}
\begin{tikzpicture}[>=Stealth,
  res/.style={draw=gcol!80!black,fill=gcol!7,rounded corners=2pt,align=center,text width=3.0cm,minimum height=1.42cm,inner sep=2pt,
    font=\scriptsize,execute at begin node=\setlength{\baselineskip}{8.6pt}},
  adp/.style={draw=black!55,fill=white,rounded corners=2pt,align=center,text width=3.0cm,minimum height=1.42cm,inner sep=2pt,
    font=\scriptsize,execute at begin node=\setlength{\baselineskip}{8.6pt}},
  pro/.style={draw=gcol!80!black,fill=gcol!7,rounded corners=2pt,align=center,text width=3.0cm,minimum height=0.9cm,inner sep=2pt,
    font=\tiny,execute at begin node=\setlength{\baselineskip}{6.9pt}},
  nog/.style={draw=ngcol,fill=ngcol!7,rounded corners=2pt,align=center,text width=3.0cm,minimum height=0.9cm,inner sep=2pt,
    font=\tiny,execute at begin node=\setlength{\baselineskip}{6.9pt}},
  opb/.style={draw=mgray,fill=lgray!45,rounded corners=2pt,align=center,text width=3.0cm,minimum height=0.9cm,inner sep=2pt,
    font=\tiny,execute at begin node=\setlength{\baselineskip}{6.9pt}},
  meet/.style={rounded corners=2pt,align=center,text width=2.8cm,minimum height=1.0cm,inner sep=2pt,anchor=north,
    font=\tiny,execute at begin node=\setlength{\baselineskip}{6.9pt}},
  est/.style={->,very thick,gcol},
  con/.style={->,very thick,tcol},
  cnd/.style={->,very thick,ocol,dashed},
  obs/.style={->,very thick,scol,densely dotted},
  opn/.style={->,very thick,mgray,dotted},
  pdown/.style={->,thick,gcol!80!black},
  ndown/.style={->,thick,ngcol,densely dotted},
  odown/.style={->,thick,mgray,dotted},
  st/.style={font=\tiny,above,inner sep=1.5pt,align=center},
  tg/.style={font=\tiny\itshape,below,inner sep=1.5pt,align=center,text=black!75},
  hd/.style={anchor=south west,font=\footnotesize\bfseries,inner sep=0pt}]
\matrix (m) [matrix of nodes,nodes in empty cells=false,column sep=1.0cm,row sep=0.18cm,cells={anchor=center}]{
% chain 1: spacetime
|[res]| {\textbf{Clifford data}\\ $\Cl(3,3)\cong M_8(\R)$ (inherited); $\Z_2^3$ grading; the roots $\cK$, $G$ (Prop.~\labelcref{prop:twoK})}
& |[res]| {\textbf{Lorentz structure}\\ $\pm J$ selected under stated hypotheses (Thm.~\labelcref{thm:i}); real Dirac $=$ complex Dirac (Thm.~\labelcref{thm:realification})}
& |[res]| {\textbf{observable four-plane}\\ the selector is $J$ (Thm.~\labelcref{thm:selector}); the light cone as twisted squares (Thm.~\labelcref{thm:cone})}
& |[res]| {\textbf{causal evolution}\\ scalar matter: Klein--Gordon field on $(M^4,h)$; unitary and microcausal for flat $h$ (Cor.~\labelcref{cor:weak})} \\
|[pro]| {also: charge conjugation is the surplus plane; time reversal squares to $-1$ (\S\labelcref{sec:discrete})}
& |[nog]| {no-go: unreduced $\R^{3,3}$ is ill-posed, with no time orientation (Props.~\labelcref{prop:toy,prop:orient})}
& |[nog]| {no-go: no unconstrained invariant adjoint scalar selects the type (Thms.~\labelcref{thm:degree,thm:flat})}
& |[opb]| {open: interacting fields; a local clock symmetry; a universal $\tau_U$ (\S\labelcref{sec:universaltau})} \\[0.3cm]
% chain 2: electromagnetism
|[res]| {\textbf{field-strength reading}\\ $\R^{3,3}=\Lambda^2\R^{1,3}$ under $\lG$, $G$ the Hodge star (Thm.~\labelcref{thm:emfield})}
& |[res]| {\textbf{charge and gauge geometry}\\ electric charge among the compact generators at aligned vacua (Thm.~\labelcref{p3:thm:photon})}
& |[res]| {\textbf{Maxwell dynamics}\\ covariance fixes Maxwell's equations (Cor.~\labelcref{cor:maxwellform}); Maxwell term adopted}
& |[res]| {\textbf{physical photon}\\ aligned vacua: photon massless with $W$ and $Z$ (Prop.~\labelcref{p3:prop:neutral})} \\
|[pro]| {also: spin one appears twice, joined by the Klein correspondence (Thm.~\labelcref{thm:klein})}
& |[nog]| {no-go: the normal charge is half the fermion number, anomalous (Rems.~\labelcref{rem:photonn3,p3:rem:anomaly})}
& |[nog]| {no-go: the normal $U(1)$ is unseen by reduced spinless matter and a ghost (Thm.~\labelcref{thm:ghostvec})}
& |[opb]| {open: the electroweak breaking mechanism; quantisation} \\[0.3cm]
% chain 3: matter and interactions
|[res]| {\textbf{tower and module}\\ $\Cl(3n,3n)$ on $(\R^8)^{\otimes n}$ and its dual pair, inherited (\S\labelcref{sec:tower})}
& |[res]| {\textbf{chiral content}\\ $P_*$, rank $64$: one Pati--Salam family with $\nu_R$ (Thm.~\labelcref{thm:PS}, Cor.~\labelcref{cor:PSfamily})}
& |[res]| {\textbf{anomaly-free gauge action}\\ Pati--Salam algebra, dimension $21$, inside the faithful $23$ (Thm.~\labelcref{thm:PSalg}); action \eqref{eq:Saction} adopted}
& |[adp]| {\textbf{interactions}\\ Yukawa couplings and potential adopted} \\
|[pro]| {also: centraliser $\so{3n-1,3n-3}$ of the interior Lorentz action, nested (Thm.~\labelcref{thm:centraliser})}
& |[nog]| {no-go: chirality imposed, not selected (Prop.~\labelcref{prop:chiralpreserved}); no $\su3$ on the axial members (Cor.~\labelcref{cor:nocolour})}
& |[nog]| {no-go: the central $\uu1$s of $P_*$ cannot be gauged on its matter (Prop.~\labelcref{prop:PSanom})}
& |[opb]| {open: the origin of the couplings and of the scales} \\[0.3cm]
% chain 4: energy and mass
|[res]| {\textbf{geometric operators}\\ roots $\cK$, $G$ (Prop.~\labelcref{prop:twoK}); Schur fixes $c\cdot\id$ (Thm.~\labelcref{thm:onec})}
& |[res]| {\textbf{interior and field dynamics}\\ binding in neutral signature (Thm.~\labelcref{thm:neutral}); the closure a pendulum (Thm.~\labelcref{thm:pendulum})}
& |[res]| {\textbf{energy and mass kernels}\\ $E_6^2=E_0^2+|\mathbf S|^2$ (Thm.~\labelcref{en:min}); a kernel space of dimension $272$ (Thm.~\labelcref{thm:kernelspace})}
& |[res]| {\textbf{stable physical states}\\ Standard-Model orbit minimal for open couplings (Thms.~\labelcref{thm:select0,thm:compat})} \\
|[pro]| {also: the Dirac mass form is $\so{3,3}$-invariant; a mass breaks no spacetime symmetry (Cor.~\labelcref{cor:fullmass})}
& |[nog]| {no-go: the interior energy has signature $(6,6)$, unbounded below (\S\labelcref{sec:esmap:distinctions})}
& |[nog]| {no-go: no mass from the $b$-field (Prop.~\labelcref{bf:alone}) or from non-orthogonal $\bE$, $\bB$ (\S\labelcref{sec:orth:nomass})}
& |[opb]| {open: global minimality on $\mathcal M$ (a conjecture); quantum stability} \\[0.3cm]
% chain 5: gravity
|[res]| {\textbf{curved geometry}\\ para-Hermitian structure; a normal $U(1)$ model, exact Einstein tensor (Thm.~\labelcref{thm:ein6})}
& |[res]| {\textbf{curvature identity}\\ curvature operator on the frame; Bianchi as $\Tr(G\cR)=0$ (Prop.~\labelcref{prop:curvop})}
& |[res]| {\textbf{variational field equations}\\ adopted Palatini--Holst action: Einstein's equations for $\sin2\theta\neq0$ (Prop.~\labelcref{prop:holst})}
& |[adp]| {\textbf{healthy gravitational dynamics}\\ not derived: matter coupling, modes, low-energy limit (\S\labelcref{sec:gravlimits})} \\
|[nog]| {no-go: transport cannot preserve the paired-mode projector where $F_N\neq0$ (Prop.~\labelcref{prop:obstruct})}
& |[pro]| {also: the double dual gives the Einstein tensor (Thms.~\labelcref{thm:phase,thm:helmholtz})}
& |[nog]| {no-go: normal-plane vectors carry negative energy relative to Einstein--Hilbert (Cor.~\labelcref{cor:ein6})}
& |[opb]| {open: a gravitational action from the arena, with its matter as source} \\
};
% chain headings, above the tallest box of each row
\path let \p0=(m-1-1.north west),\p1=(m-1-1.north),\p2=(m-1-2.north),\p3=(m-1-3.north),\p4=(m-1-4.north),
  \n1={max(\y1,\y2,\y3,\y4)} in node[hd] at (\x0,\n1+0.12cm) {1\enspace Spacetime: from Clifford data to causal evolution};
\path let \p0=(m-3-1.north west),\p1=(m-3-1.north),\p2=(m-3-2.north),\p3=(m-3-3.north),\p4=(m-3-4.north),
  \n1={max(\y1,\y2,\y3,\y4)} in node[hd] at (\x0,\n1+0.12cm) {2\enspace Electromagnetism: from the field-strength reading to the physical photon};
\path let \p0=(m-5-1.north west),\p1=(m-5-1.north),\p2=(m-5-2.north),\p3=(m-5-3.north),\p4=(m-5-4.north),
  \n1={max(\y1,\y2,\y3,\y4)} in node[hd] at (\x0,\n1+0.12cm) {3\enspace Matter and interactions: from the tower to interactions};
\path let \p0=(m-7-1.north west),\p1=(m-7-1.north),\p2=(m-7-2.north),\p3=(m-7-3.north),\p4=(m-7-4.north),
  \n1={max(\y1,\y2,\y3,\y4)} in node[hd] at (\x0,\n1+0.12cm) {4\enspace Energy and mass: from geometric operators to stable physical states};
\path let \p0=(m-9-1.north west),\p1=(m-9-1.north),\p2=(m-9-2.north),\p3=(m-9-3.north),\p4=(m-9-4.north),
  \n1={max(\y1,\y2,\y3,\y4)} in node[hd] at (\x0,\n1+0.12cm) {5\enspace Gravity: from curved geometry to gravitational dynamics};
% links along the chains
\draw[cnd] (m-1-1) -- node[st]{estab.\\ cond.} node[tg]{clock\\ chosen} (m-1-2);
\draw[cnd] (m-1-2) -- node[st]{cond.} node[tg]{reduction\\ imposed} (m-1-3);
\draw[opn] (m-1-3) -- node[st]{cond.\\ open} node[tg]{gauge\\ chosen} (m-1-4);
\draw[cnd] (m-3-1) -- node[st]{cond.} node[tg]{no map\\ joins} (m-3-2);
\draw[cnd] (m-3-2) -- node[st]{cond.} node[tg]{term\\ adopted} (m-3-3);
\draw[opn] (m-3-3) -- node[st]{open} node[tg]{no EW\\ breaking} (m-3-4);
\draw[con] (m-5-1) -- node[st]{constr.} node[tg]{$P_*$\\ chosen} (m-5-2);
\draw[cnd] (m-5-2) -- node[st]{cond.} node[tg]{action\\ adopted} (m-5-3);
\draw[opn] (m-5-3) -- node[st]{open} node[tg]{couplings\\ adopted} (m-5-4);
\draw[con] (m-7-1) -- node[st]{constr.} node[tg]{premises\\ chosen} (m-7-2);
\draw[cnd] (m-7-2) -- node[st]{estab.\\ cond.} node[tg]{mass an\\ input} (m-7-3);
\draw[opn] (m-7-3) -- node[st]{cond.\\ open} node[tg]{potential\\ adopted} (m-7-4);
\draw[est] (m-9-1) -- node[st]{estab.} node[tg]{given\\ metrics} (m-9-2);
\draw[opn] (m-9-2) -- node[st]{open} node[tg]{action\\ adopted} (m-9-3);
\draw[obs] (m-9-3) -- node[st]{obstr.\\ open} node[tg]{normal\\ plane} (m-9-4);
% further results, no-go theorems and open problems below each stage
\draw[pdown] (m-1-1) -- (m-2-1); \draw[ndown] (m-1-2) -- (m-2-2); \draw[ndown] (m-1-3) -- (m-2-3); \draw[odown] (m-1-4) -- (m-2-4);
\draw[pdown] (m-3-1) -- (m-4-1); \draw[ndown] (m-3-2) -- (m-4-2); \draw[ndown] (m-3-3) -- (m-4-3); \draw[odown] (m-3-4) -- (m-4-4);
\draw[pdown] (m-5-1) -- (m-6-1); \draw[ndown] (m-5-2) -- (m-6-2); \draw[ndown] (m-5-3) -- (m-6-3); \draw[odown] (m-5-4) -- (m-6-4);
\draw[pdown] (m-7-1) -- (m-8-1); \draw[ndown] (m-7-2) -- (m-8-2); \draw[ndown] (m-7-3) -- (m-8-3); \draw[odown] (m-7-4) -- (m-8-4);
\draw[ndown] (m-9-1) -- (m-10-1); \draw[pdown] (m-9-2) -- (m-10-2); \draw[ndown] (m-9-3) -- (m-10-3); \draw[odown] (m-9-4) -- (m-10-4);
% where the chains meet
\node[hd,anchor=north west] (wm) at ([yshift=-0.32cm]m.south west) {Where the chains meet};
\node[meet,draw=gcol!80!black,fill=gcol!7,anchor=north west] (w1) at ([yshift=-0.12cm]wm.south west)
  {\textbf{one complex unit}\\ chains 1, 3: the imaginary unit $J$ of Dirac theory selects the four-plane (Thms.~\labelcref{thm:i,thm:selector})};
\node[meet,draw=gcol!80!black,fill=gcol!7,anchor=north west] (w2) at ([xshift=0.28cm]w1.north east)
  {\textbf{one binding root}\\ chains 2, 4: $G$ binds the interior dynamics and is the Hodge star (Thms.~\labelcref{thm:neutral,thm:emfield})};
\node[meet,draw=tcol,fill=tcol!6,anchor=north west] (w3) at ([xshift=0.28cm]w2.north east)
  {\textbf{one joint realisation}\\ chains 2, 3: two-flavour spinor electrodynamics at the second rung (Thms.~\labelcref{thm:mixeddirac,thm:coupled})};
\node[meet,draw=ngcol,fill=ngcol!7,anchor=north west] (w4) at ([xshift=0.28cm]w3.north east)
  {\textbf{two readings}\\ chains 1, 2: $\lint$ and $\lG$ are not conjugate; no nonzero equivariant linear map (Props.~\labelcref{prop:notconj,prop:emexclusive})};
\node[meet,draw=ngcol,fill=ngcol!7,anchor=north west] (w5) at ([xshift=0.28cm]w4.north east)
  {\textbf{matter and binding}\\ chains 3, 4: the centraliser of $G_3$ meets $\lint$ in $0$ (Thms.~\labelcref{thm:nocommon,thm:PSbinding})};
\node[anchor=north west,inner sep=0pt,font=\tiny,text width=15.45cm,align=left] (ch) at ([yshift=-0.15cm]w1.south west)
  {\emph{Choices left free by the algebra} (\cref{tab:choices}): clock, projector, Lorentz action, hypercharge direction, binding operator,
  scales and inertias, action, four-plane and positivity, families.};
% legend
\begin{scope}[shift={([yshift=-0.32cm]ch.south west)},font=\scriptsize]
\draw[est] (0,0) -- (0.55,0); \node[anchor=west] at (0.6,0) {established};
\draw[con] (2.35,0) -- (2.9,0); \node[anchor=west] at (2.95,0) {constructed};
\draw[cnd] (4.75,0) -- (5.3,0); \node[anchor=west] at (5.35,0) {conditional};
\draw[obs] (7.1,0) -- (7.65,0); \node[anchor=west] at (7.7,0) {obstructed};
\draw[opn] (9.4,0) -- (9.95,0); \node[anchor=west] at (10.0,0) {open};
\node[draw=gcol!80!black,fill=gcol!7,rounded corners=1pt,font=\tiny,inner sep=2pt] at (11.3,0) {proved};
\node[draw=ngcol,fill=ngcol!7,rounded corners=1pt,font=\tiny,inner sep=2pt] at (12.75,0) {no-go theorem};
\node[draw=mgray,fill=lgray!45,rounded corners=1pt,font=\tiny,inner sep=2pt] at (14.05,0) {open};
\node[font=\tiny\itshape,text=black!75] at (15.2,0) {choice made};
\end{scope}
\end{tikzpicture}
\caption{The research chains of \cref{sec:chains} in detail. Green: results proved at each stage, under the hypotheses stated where they
are proved, and in the bottom row results that join the chains; violet: no-go theorems, proved obstructions that close a route; grey: open.
Each arrow carries its weakest status (both where \cref{tab:chains} gives two) and, below it, the choice on which the link rests. A summary,
not a derivation; no chain is complete.}
\label{fig:chainsdetail}
\end{figure}

{\footnotesize
\begin{longtable}{@{}>{\raggedright\arraybackslash}p{3.3cm}>{\raggedright\arraybackslash}p{2.5cm}>{\raggedright\arraybackslash}p{8.9cm}@{}}
\caption{The links of the five research chains with their statuses and the places where each is proved or qualified.}\label{tab:chains}\\
\toprule
link & status & evidence and limits\\
\midrule
\endfirsthead
\multicolumn{3}{@{}l}{\emph{\cref{tab:chains}, continued}}\\
\toprule
link & status & evidence and limits\\
\midrule
\endhead
\bottomrule
\endlastfoot
\multicolumn{3}{@{}l}{\emph{Spacetime}}\\
Clifford data $\to$ Lorentz structure & established (existence); conditional (selection) & the interior $\so{1,3}$ and $\lG$ exist (\cref{ch:ch01_clifford,ch:ch02_lorentz,ch:ch03_readings}); choosing one needs a clock and stated hypotheses, and all unit clocks are equivalent under rotations (\cref{rem:clockcov}); the imaginary unit is unique only under the hypotheses of \cref{thm:i}\\
Lorentz structure $\to$ observable four-plane & conditional & $E_{\hat n}$ is selected by a clock; the clock reduction is \cref{assum:clockred}, imposed; no dynamics of an unconstrained $SO(3,3)$-invariant adjoint scalar selects the four-plane (\cref{thm:degree,thm:flat,prop:kineticB})\\
four-plane $\to$ causal evolution & conditional (scalar matter); open (interacting) & weak emergence for scalar matter under the chosen gauge principle (\cref{cor:weak}); constraint preservation shown for a test worldline only (\cref{thm:wlgauge}); no local clock symmetry is constructed (\cref{sec:localclock}); a universal evolution parameter, prior art in the theory of Horwitz and Piron, is an open hypothesis in this arena (\cref{sec:universaltau})\\
\addlinespace
\multicolumn{3}{@{}l}{\emph{Electromagnetism}}\\
field-strength reading $\to$ charge and gauge geometry & established (each part); conditional (the link); obstructed (normal charge as photon) & $\R^{3,3}=\Lambda^2\R^{1,3}$ under $\lG$ with $G$ the Hodge star (\cref{thm:emfield}), a representation theorem; the electric charge is identified among the compact generators at the aligned vacua (\cref{p3:thm:photon}); no map or common action joins the two, so the link is conditional; the normal charge is half the fermion number on $V_*$ and anomalous with $\su2_L$ (\cref{rem:photonn3,p3:rem:anomaly}), and the normal $U(1)$ of the curved model fails on three further counts (\cref{sec:transport,thm:ghostvec})\\
charge $\to$ Maxwell dynamics & conditional & covariance fixes the form of Maxwell's equations (\cref{cor:maxwellform}); the Maxwell term is adopted, not derived\\
Maxwell dynamics $\to$ physical photon & open & the aligned vacuum is massless for the photon together with $W$ and $Z$ and is electroweak symmetric; the adopted potential opposes the neutral deformation (\cref{p3:prop:neutral}); the electroweak breaking mechanism is missing; no quantisation\\
\addlinespace
\multicolumn{3}{@{}l}{\emph{Matter and interactions}}\\
tower and module $\to$ chiral content & constructed & $P_*$ is chosen among projectors; it carries one Pati--Salam family with $\nu_R$ (\cref{thm:PS,cor:PSfamily}); chirality is imposed by the choice, not selected by dynamics (\cref{prop:chiralpreserved}); the diagonal Lorentz action gives vector-like matter instead (\cref{thm:diagvectorlike})\\
chiral content $\to$ anomaly-free gauge action & conditional & the Standard-Model algebra is anomaly free on $V_*$; the two central $\uu1$s of $P_*$ cannot be gauged on it (\cref{prop:PSanom}); the gauge action \eqref{eq:Saction} is adopted\\
gauge action $\to$ interactions & open & Yukawa couplings and the scalar potential are adopted; no coupling or scale relation is derived\\
\addlinespace
\multicolumn{3}{@{}l}{\emph{Energy and mass}}\\
geometric operators $\to$ interior and field dynamics & constructed & the closure is unique given its premises and its free motion is a pendulum (\cref{thm:closureunique,thm:pendulum}); bounded orbits do not make the energy bounded below, which has signature $(6,6)$ (\cref{sec:esmap:distinctions})\\
dynamics $\to$ energy and mass kernels & established (classification); conditional (link to dynamics) & the $272$-dimensional kernel space and its complex structure $\widetilde{\mathcal C}$ (\cref{thm:kernelspace}); a mass in a Dirac bilinear is an input, not generated; $E_6^2=E_0^2+|\mathbf S|^2$ is an identity, not a process (\cref{sec:resynthesis})\\
kernels $\to$ stable physical states & conditional (selection); open (stability) & the Standard-Model orbit is the minimum of the adopted potential on a stated open set of couplings (\cref{thm:select0,thm:compat}); global minimality on all of $\mathcal M$ is a conjecture supported by numerical searches (\cref{prop:randomstart}), within the model assumptions stated in \cref{ch:ch07_vacuum}; quantum stability is open\\
\addlinespace
\multicolumn{3}{@{}l}{\emph{Gravity}}\\
curved geometry $\to$ curvature identity & established & the ambient Einstein tensor of the model and the double dual of the field-strength reading (\cref{thm:ein6,thm:phase,thm:helmholtz})\\
curvature identity $\to$ variational field equations & open & no gravitational action is derived; the Palatini--Holst action is adopted (\cref{prop:holst,tab:gravlimits})\\
field equations $\to$ healthy gravitational dynamics & obstructed (normal-plane route); open (otherwise) & relative to the ambient Einstein--Hilbert action and the stated ansatz, every normal-plane vector gravitates with negative energy, and the sign is not a convention (\cref{thm:ghostvec,rem:signconv}); matter coupling, degrees of freedom and low-energy limit are not derived (\cref{sec:gravlimits})\\
\end{longtable}}

No chain is complete. Two features recur in all five. Every chain passes through a choice that the algebra does not make, the clock,
the projector, the Lorentz action, an action or a potential (\cref{tab:choices}); and every chain stops before dynamics that would be derived
rather than adopted. The chains also do not share their data: the spacetime chain uses $\lint$ after a clock, the electromagnetic chain uses
$\lG$, and the two are not conjugate (\cref{prop:notconj,prop:emexclusive}). \Cref{sec:incompatible} lists the places where this matters.

\section{Sectors, representations and projectors}
\label{sec:sectors}

\Cref{tab:sectors} records for each sector the data that must agree before two sectors can be combined. It is the manuscript form of
the compatibility map kept by the author.

\begin{table}[p]
\centering
\scriptsize
\begin{tabular}{@{}>{\raggedright\arraybackslash}p{1.9cm}>{\raggedright\arraybackslash}p{2.3cm}>{\raggedright\arraybackslash}p{1.8cm}>{\raggedright\arraybackslash}p{2.0cm}>{\raggedright\arraybackslash}p{2.5cm}>{\raggedright\arraybackslash}p{3.2cm}@{}}
\toprule
sector & algebra and module & Lorentz action & complex structure; projector & form; action & status\\
\midrule
displacement reading (\cref{ch:ch03_readings}) & $V=\R^{3,3}$, vector module & $\lint(\hat n)$ after a clock & none; $\pr_{\hat n}$ & $\eta$; free worldline, shifted interior Lagrangian & kinematics established; interior energy indefinite\\
field-strength reading (\cref{ch:ch03_readings}) & $V\cong\Lambda^2E_{\hat n}^*$ & $\lG$, $(1,0)\oplus(0,1)$ & $G$ acts as $i$ on $z$; none & $\eta$, $q=z\cdot z$; Maxwell (inherited) & established; no nonzero Lorentz-equivariant linear intertwiner to the displacement reading (\cref{r1:thm:type})\\
field octants (\cref{ch:ch03_readings}) & $S\cong\Lambda^*(\mathsf L)=\R^8$ & $\rho_{\rm oct}$, with odd $Q_i$ & none chosen (the duality lift of \cref{ch:ch03_readings} exists) & $T$, $T\Gamma$; Maxwell on the octants & fixed; representation-specific\\
Dirac field, one constituent (\cref{ch:ch02_lorentz}) & $S=\R^8$, Grassmann field & $\rho_{\hat n}$, clock-dependent & $J(\hat n)=T\hat n$; none & mass form $T$; real Dirac Lagrangian & equivalent to complex Dirac theory (\cref{thm:realification}); $m$ an input\\
tower, $n\ge2$ (\cref{ch:ch05_tower}) & $\Cl(3n,3n)$ on $(\R^8)^{\otimes n}$ & per constituent or diagonal & selected unit of the two-particle space; Fock-number projectors & graded or ungraded rest operators; no dynamics beyond free kernels & algebra proved; particle reading conditional\\
chiral block, $n=3$ (\cref{ch:ch06_chiral}) & $\Cl(9,9)$ on $\R^{512}$; $V_*=\operatorname{im}P_*$, rank $64$ & interior $\so{1,3}$ & $J$ (14 internal generators); $P_*$ & $K^\mu=P_*\gamma^0\gamma^\mu P_*$; action \eqref{eq:Saction} adopted & content proved (\cref{thm:PS}); preserved by compact Pati--Salam only, not by $\so{4,6}$ (z3-9)\\
family block (\cref{ch:ch06_chiral}) & $\operatorname{im}P_{\rm fam}$, rank $64$ & interior $\so{1,3}$ & $P_{\rm fam}$, not $P_*$ & as above & proved; a different orbit from $P_*$\\
mass kernels (\cref{ch:ch07_vacuum}) & $\mathcal M$, real dimension $272$ & invariant by definition & $\widetilde{\mathcal C}$ on $\mathcal M$; none & $\Tr M^{\!\top}N$; adopted potential, leading-log selection & classification proved (\cref{thm:kernelspace}); selection conditional; global minimality conjectural\\
compact vacuum families, photon (\cref{ch:ch08_photon}) & projector fields on the family; aligned points $\Pi_{\pm A},\Pi_{\pm W}$ & -- & --; $\Pi$, not $P_*$ & trace form; adopted $V_{\mathfrak h}$ and Maxwell term & charge identified; normal charge anomalous; electroweak breaking mechanism missing\\
observational projection (\cref{ch:ch02_lorentz,ch:ch12_causal}) & $\pr_{\hat n}$ on $V$; reduced shell & $\lint(\hat n)$ & --; $\pr_{\hat n}$ & $\eta|_{E_{\hat n}}$; surplus constraints & imposed (\cref{assum:clockred}); worldline scope\\
interior dynamics (\cref{ch:ch04_rigidbody}) & frame and rotor on $\so{3,3}$; closure $\mathfrak s$ & $\lG$ & $G$; none & Lie--Poisson; inertia $\mathbb I$ & bounded orbits proved; Hamiltonian not bounded below\\
$b$-field (\cref{ch:ch09_bfield}) & connections commuting or anticommuting with $G$; copy rotations & per constituent & $G$ per constituent; none & split virial identity & proved in the stated classes; rates not derived\\
curved arena (\cref{ch:ch11_curvature}) & para-complex $\cK$; normal $U(1)$ model & local $\lint$ & $\cK$ (para-complex); $\Pi(x)$ transported & $g_6$; Einstein-tensor contractions & identities established; no gravitational action\\
causal structure (\cref{ch:ch12_causal}) & wave operator on $\R^{3,3}$; selector & -- & selector $=J$; type selection & $\eta$; worldline gauge & conditional; worldline scope\\
\bottomrule
\end{tabular}
\caption{The sectors of the manuscript and the data that must agree before two of them can be combined.}
\label{tab:sectors}
\end{table}

\paragraph{Projectors kept apart.} Six families of projectors occur, and no statement transfers from one to another without an explicit map.
\begin{enumerate}[label=(\roman*),leftmargin=2.2em]
\item $P_*$, the chiral matter block of \cref{sec:pstar}, a choice among the projectors that commute with the Lorentz action and $J$.
\item $P_{\rm fam}$, the family block of \cref{ch:ch06_chiral}: of the same rank, on a different orbit, with a different effective algebra.
\item The auxiliary projectors $P_{ij}$, $P_{L,R}$, $P_\nu$ on the module (\cref{ch:ch06_chiral,ch:ch07_vacuum}).
\item The vacuum projectors $\Pi$ of the compact families, with the aligned points $\Pi_{\pm A}$, $\Pi_{\pm W}$ (\cref{ch:ch08_photon}); their overlaps with $P_L$ and
$P_{\rm fam}$ enter the adopted potential, which encodes the target of the selection in its couplings.
\item The observational projection $\pr_{\hat n}$ onto the four-plane of a clock (\cref{ch:ch02_lorentz,ch:ch12_causal}), which acts on the vector space $\R^6$.
\item The transported projector field $\Pi(x)$ of \cref{sec:transport}, which acts on the spinor module $\R^8$. It is not the curved counterpart of
$\pr_{\hat n}$: \cref{ch:ch11_curvature} shows that the two behave differently under transport.
\end{enumerate}
A transported projector keeps its rank and idempotence, but a projector fixed once, such as $P_*$, is not invariant under clock rotations
unless it is transported with them (\cref{rem:clockcov}(b)).

\section{Incompatible readings and missing bridges}
\label{sec:incompatible}

The sectors of \cref{tab:sectors} are realised on one algebra, but several of them are not realised together. The following
incompatibilities are proved, and each is part of the result.
\begin{enumerate}[label=(\arabic*),leftmargin=2.2em]
\item \emph{Displacement and field strength.} The two readings of the same six directions share $\eta$ and the Clifford action, but their
Lorentz algebras $\lint$ and $\lG$ are not conjugate and no nonzero Lorentz-equivariant linear intertwiner joins the specified modules (\cref{prop:notconj,prop:emexclusive,r1:thm:type}). They
meet only as base and fibre, which is a reading.
\item \emph{Spin one-half and the octants.} $\rho_{\hat n}$ and $\rho_{\rm oct}$ are two Lorentz representations on one module; the shared
operator $T$ does not identify them.
\item \emph{Interior energy and field energy.} The interior energy has signature $(6,6)$ and the field energy $E_6$ is positive; no map makes
one the other, and the two forms lie in the Lorentz class and the duality class of \cref{thm:inertiareq}, which meet in $0$.
\item \emph{Matter and the interior dynamics.} The chiral block is built on $\lint$, the binding and the closure on $\lG$. At three
constituents the centraliser of the binding root $G_3$ meets the interior Lorentz algebra of the first constituent in $0$ (\cref{thm:nocommon}), and the part of
the Pati--Salam algebra compatible with the binding is $\so4_\Delta$ (\cref{thm:PSbinding}).
\item \emph{Six-dimensional kinematics and the reduced shell.} Uphill decays exist without a common clock and disappear only because the
reduction is imposed (\cref{e7:prop:uphill}).
\item \emph{$P_*$ and $\so{4,6}$.} The block $V_*$ is not a module of the $\so{4,6}$ of Krasnov's construction; only the compact part preserves it
(\cref{rem:krasnovpointer}, z3-9). No intertwiner has been constructed.
\item \emph{The normal $U(1)$ and the photon.} The normal connection of the curved model is invisible to spinless matter in the zero-normal-momentum
reduced sector, its adapted generator
does not commute with the physical generators, relative to the ambient Einstein--Hilbert action it gravitates with the wrong sign, and on the chiral block its charge is anomalous
(\cref{sec:transport,thm:ghostvec,rem:photonn3}).
\end{enumerate}

\paragraph{Joint realisations.} One limited joint realisation is constructed. At the mixed second rung the univectorial potential, a Dirac
block and the frame field are coupled in one gauge theory, two-flavour spinor electrodynamics (\cref{thm:mixeddirac,thm:coupled}); the price is
the internal algebra, whose centraliser shrinks to the rotation of one surplus pair (\cref{prop:mixedcost}), and the block is not the chiral
block of \cref{sec:pstar}. No joint realisation of the chiral block, the field-strength reading, the interior dynamics and the curved arena
exists in this manuscript.

\paragraph{Missing bridges.} What would connect the separate realisations is, in each case, a construction that is not made:
a common Lorentz action for the two readings, or a proof that none exists at an acceptable price beyond the five candidates tested; an explicit map between the octants and the spin module that intertwines a stated dynamics; a derivation of the clock reduction
for interacting fields; an action relating metric, pairing, connection and projector; an origin of the adopted potentials; and an
intertwiner, if one exists, between $V_*$ and Krasnov's even forms. A common local action for a compatible subset of the sectors is a
possible future milestone. It is not a result here, and it is not a prerequisite for the mathematical statements of the earlier chapters.

\section{The rest-energy hypothesis, assessed}
\label{sec:resynthesis}
\label{s5:sec:synthesis}

Part of the programme was motivated by a reading of $E=mc^2$ proposed by the author and stated in the introduction
(\cref{s0:sec:question}): the energy is one undivided, six-dimensional entity; the mass emerges through an internal differentiation of it;
$c^2$ is the symmetric part of the relation; $m$ is what breaks the symmetry; and $m$ is generated from the $b$-field of Yang and Mills. The
relation $E=mc^2$ itself is inherited. What can be assessed is how the arena realises its ingredients and which realisations are forced. The
three clauses about the energy and $c^2$ are held against the results first; the two clauses about breaking and generation follow.
\label{sec:w:whatthearenasaysaboutemc}
\label{sec:energy-readings}


The results of \cref{sec:energy} can be held against three clauses that are often attached to $E=mc^2$: that the energy is one undivided,
six-dimensional quantity; that mass arises from an internal differentiation of it; and that $c^2$ is the symmetric part of the relation.

\paragraph{The energy is one undivided, six-dimensional quantity.} This holds in the field-strength reading. Among quadratic forms on the six
frame components exactly one is invariant under the diagonal rotations and under duality, and it is positive: $E_6=\tfrac12(|\bB|^2+|\bE|^2)$
(\cref{en:unique}, inherited, restating \cref{thm:inertiareq}(iii)). The transformations that preserve it and commute with $G$ form $U(3)$; its full
linear symmetry group is $O(6)$. The energy is unitary in this sense, but only as a pointwise quadratic form: $U(3)$ is not a symmetry of Maxwell's equations on the frame nor of the split-form interior (\cref{en:u3}). In
the closure the undivided energy of the momentum field, $E_6(\pi)=\tfrac12|p|^2$, is a Casimir, conserved by every Hamiltonian together with
$E_0(\pi)$ and $\mathbf S(\pi)$ (\cref{en:closure}, restating \cref{prop:LP}); the energies of the frame field are not conserved
(\cref{en:interior,en:closure}). The other reading of the clause, three temporal energies forming a vector in $V^+$, is obstructed: splitting that
vector lowers the mass (\cref{en:temporal}).

\paragraph{Mass arises from an internal differentiation of that quantity.} This holds as an identity, not as a process. With no interaction and no
second field the undivided energy of a field point splits as $E_6^2=E_0^2+|\mathbf S|^2$, $E_0=\tfrac12|z\cdot z|$ (\cref{en:split}), in SI units
$u^2=c^2|\mathbf g|^2+u_0^2$, the field-point form of $E^2=(pc)^2+(mc^2)^2$ (\cref{en:SI}). The rest part is the complete invariant of Lorentz
transformations and duality on nonzero fields, the zero field being a separate orbit (\cref{en:orbits}), the Lorentz scalar of the stress tensor whose
$00$ component is $E_6$ (\cref{en:rainich}), and the infimum of $E_6$ over boosts, attained exactly at vanishing Poynting vector for a non-null field and
not attained for a nonzero null field, which has no rest frame (\cref{en:min}). All of this is inherited classical electrodynamics in the arena's
variables. The arena adds which form plays which role: $E_6$ belongs to the duality class of \cref{thm:inertiareq} and $E_0$, the modulus of the two
Lorentz-class invariants, is built from the Lorentz class, and reading $E=mc^2$ as the locus where the two coincide, the rest frame, is a reading. The
differentiation does not happen in time: in the closure the split of the momentum field is fixed by the initial data, all three terms being
Casimirs (\cref{en:closure}). Nor does a whole get its mass from interference: two disjoint counter-propagating null pulses have $q\equiv0$ everywhere
and $M=U>0$ (\cref{en:whole}).

\paragraph{$c^2$ is the symmetric part.} This holds in one precise form. Schur's lemma fixes the isotropic form $C=c\cdot\id$ of the conversion
between the temporal and the spatial triple, not the value of $c$ (\cref{thm:onec}); $c$ is a conversion, not a symmetry, and its weight in $E_6$ is a unit convention
(\cref{en:onec}). The form in which the clause is exact is the decomposition of the shift into the two summands of \cref{thm:EAdichotomy},
\[
\End(V)=\so{3,3}\oplus\{Y:Y^\dagger=Y\},\qquad \kappa G\in\{Y:Y^\dagger=Y\},\qquad (\kappa G)^2=-\frac{c^2}{\ell^2}\,\id .
\]
The $\eta$-self-adjoint summand is the one in which factorisation holds, the one that binds in neutral signature (\cref{thm:neutral}), and the one
that carries $c$, through $\kappa=c/\ell$; the antisymmetric summand holds the connections. The decomposition and the fact that only the first summand
binds are fixed; the value $(\kappa G)^2=-c^2/\ell^2$ is constructed, a definition with $\ell$ free; and reading the first summand as ``$c^2$'' and the second, which holds the
connections, as the place of a gauge field is a reading; that a gauge generator differentiates the energy is not supported, since it selects no
amount of rest energy (\cref{en:transport}). Its limits are fixed results: duality is a symmetry of the positive form,
while the interior Lagrangian is built on $\eta$, and the two forms lie in the duality and the Lorentz class of \cref{thm:inertiareq}, which meet in $0$;
in the interior $c^2$ enters as $\kappa^2=c^2/\ell^2$ and the energy is $mc^2$ times an indefinite, unselected amplitude (\cref{en:interior}); and in the
closure what $E_0(\pi)$ breaks is the constancy of the rotor rate, not anything carried by $c$ (\cref{en:closure}).

Collected, the clauses become the following statement, which is a reading of fixed and inherited results. The undivided energy is
$E_6=\tfrac12|z|^2$, the one rotation- and duality-invariant quadratic energy of a six-component frame vector, preserved by the $U(3)$ of
transformations commuting with $G$. It
differentiates internally, with no interaction, as $E_6^2=E_0^2+|\mathbf S|^2$, into a Lorentz-invariant rest energy $E_0=\tfrac12|z\cdot z|$ and a
motion part $|\mathbf S|$. What differentiates is the arena's splitting into a temporal and a spatial triple, the real structure $\Ad_{e_1e_2e_3}$.
Within $U(3)$ the real structure is preserved by $O(3)$ and $E_0$ by $U(1)\cdot O(3)$, symmetries of pointwise forms and not of any dynamics; nothing Lorentzian is broken, and the duality
$U(1)_G$ is broken to $\Z_2$ only on the space of Lorentz orbits and only by a configuration. The constant $c$ is a conversion, not a symmetry, and
$E_0/c^2$ is a rest-mass density of a field point. The statement has three limits: a field point is not a particle (two boost factors, and the mass of
a whole is not interference); in the closure only the momentum field's energies are conserved, while the frame field's oscillate or grow; and a compact
gauge generator selects no amount of rest energy, its anti-diagonal part moving a field between null and rest at fixed $E_6$ and, at three
constituents, $\so4_\Delta$ changing the rest energies of the constituents at fixed $\tfrac12|Z|^2$ and $q_3$. Whether some dynamics selects the rest
frame of a field or the value of $E_0$ is open; nothing in this manuscript does.

\paragraph{$m$ is what breaks the symmetry.} This holds for internal symmetries and fails for spacetime ones (\cref{sec:massbreaks}).
\begin{enumerate}[label=(\roman*),leftmargin=2.2em]
\item A nonzero rest energy breaks duality: the duality $U(1)_G$ acts on the Lorentz invariant by $q\mapsto e^{2i\theta}q$, so for $E_0\neq0$ the
stabiliser of $q$, or of the Lorentz orbit it labels, is $\{\pm\id\}$ (\cref{en:doubly,en:stab}). Under duality alone a nonzero field point has trivial
stabiliser, since $e^{\pi G}F=-F$. At a null field point duality joins the little group. This is inherited, pulled back to the frame.
\item For Dirac matter the mass breaks the chiral and Pauli--G\"ursey symmetries, $\uu1_\Omega\oplus\su2\to\uu1_J$, with three broken
directions (\cref{sec:massdirac,sec:massgoldstone}); for an explicit Dirac mass these are directions of the mass orbit, not Goldstone excitations.
\item The mass breaks no spacetime symmetry, not even the six-dimensional one: the Dirac mass form is the only Grassmann mass kernel invariant
under all of $\so{3,3}$ (\cref{cor:fullmass}). Within $\so{3,3}$ it is the kinetic term that cuts the symmetry down. In this sense the clause is
inverted.
\item In the closure the rest energy of the momentum field is the amplitude of the potential of the duality angle, so a null momentum field
leaves the rate constant and a massive one makes the angle a pendulum (\cref{thm:pendulum,cor:constrate}). Reading the angle as a
pseudo-Goldstone boson of duality is interpretive: the breaking is explicit, carried by the Lorentz-class inertia.
\end{enumerate}

\paragraph{$m$ is generated from the $b$-field.} Classically this fails, and the failure is sharp (\cref{sec:bfield}).
\begin{enumerate}[label=(\roman*),leftmargin=2.2em]
\item The binding rate is never a gauge potential: for every connection $b$ the field equation factorises covariantly, $\kappa G$ transforming as a
tensor and $b$ as a connection, and without $\kappa$ a constant $b$ is pure gauge along the worldline (\cref{bf:covariant}).
\item An admissible connection is compact, by the definition of gauge admissibility after \cref{cor:noboost}. Its part commuting with $G$ splits the rest
frequency into charge-labelled lines and is removable along one worldline by a rotating frame (\cref{bf:larmor}); the screening criterion of
\cref{bf:screen} decides when a mixed connection keeps every orbit bounded; and a connection alone binds nothing (\cref{bf:alone}).
\item At three constituents the part of the Pati--Salam algebra commuting with the binding root is $\so4_\Delta\cong\su2\oplus\su2$
(\cref{thm:PSbinding}); it splits the rest frequency by weights and supplies no scale (\cref{bf:su2}).
\item In the selected vacuum the order is reversed: the Pati--Salam vector bosons acquire mass by eating Goldstone directions
(\cref{sec:goldstone,thm:nomajoron}) rather than generating the mass of anything.
\end{enumerate}
The one role left for the $b$-field is the scale. The algebra fixes the isotropic form of the conversion map, while its normalisation and the
dynamical length scale $\ell$, hence $\kappa=c/\ell$, require additional input; the identification of
$\hbar\kappa$ with a scale generated by dimensional transmutation is a conjecture (\cref{bf:transmute}), and reading the rate as a mass,
$mc^2=\hbar\kappa$, is a further hypothesis (\cref{sec:restscale}).

\begin{table}[t]
\centering
\small
\begin{tabular}{@{}>{\raggedright\arraybackslash}p{3.4cm}>{\raggedright\arraybackslash}p{6.7cm}>{\raggedright\arraybackslash}p{4.2cm}@{}}
\toprule
clause & what the arena gives & verdict\\
\midrule
energy undivided, six-dimensional & $E_6=\tfrac12(|\bB|^2+|\bE|^2)$, the unique rotation- and duality-invariant quadratic energy; $U(3)$ preserves it and commutes with $G$ (full group $O(6)$); a Casimir of the closure for the momentum field; the temporal-vector reading lowers the mass & holds in the field-strength reading (inherited); the temporal reading is obstructed\\
mass from internal differentiation & $E_6^2=E_0^2+|\mathbf S|^2$, $E_0=\inf_{\rm boosts}E_6$, attained at vanishing Poynting vector and not for a nonzero null field; in the closure the split is fixed by the initial data & holds as an identity, not as a process\\
$c^2$ the symmetric part & Schur's lemma fixes the isotropic form $c\cdot\id$ of the conversion, not the value of $c$ (\cref{thm:onec}); $\kappa G$ the $\eta$-self-adjoint summand of the shift, the only one that binds & the decomposition is fixed; the reading is interpretive\\
$m$ breaks the symmetry & breaks duality and the chiral algebra; breaks neither Lorentz nor $\mathrm{Spin}(3,3)$ & true for internal symmetries, false for spacetime ones\\
$m$ generated from the $b$-field & the connection summand splits or tears the binding and never creates it; vector bosons acquire mass by eating Goldstones & fails classically; a scale by transmutation is conjectural\\
\bottomrule
\end{tabular}
\caption{The rest-energy hypothesis against the arena (\cref{sec:resynthesis}).}
\label{tab:resynthesis}
\end{table}

Collected, the strongest statement of the hypothesis that the arena supports is the following (\cref{tab:resynthesis}). The undivided energy is
the duality-invariant energy $\tfrac12(|\bB|^2+|\bE|^2)$ of the six field components. The conversion between times and lengths is isotropic,
$c\cdot\id$, by Schur's lemma, which does not fix the value of $c$, and it enters the dynamics through the $\eta$-self-adjoint part of the shift, the binding root. The mass is the
Lorentz- and duality-invariant $\tfrac12|q|$, the infimum of the energy over boosts, attained in the rest frame of a non-null field, and $E=mc^2$ is
the rest frame, where the two coincide.
What the mass breaks is duality and, for matter, chirality, never Lorentz invariance or $\mathrm{Spin}(3,3)$. The $b$-field is the
antisymmetric part of the same shift: it can split or destabilise the binding but cannot create it. Every part of this statement is fixed,
inherited or interpretive as marked; none of it is a derivation of a mass.

\section{Duplication and the \texorpdfstring{$b$}{b}-field between manifolds}
\label{sec:dupsynthesis}
\label{s6:sec:synthesis}

The second question of the motivating hypothesis (\cref{s0:sec:question}) asked what process makes a unitary energy result in a duplication,
from $\Cl(3,3)$ to $\Cl(6,6)$, and whether the $b$-field is a deal breaker between two manifolds or more. The answers of
\cref{ch:ch05_tower,ch:ch09_bfield} are collected in \cref{tab:dupsynthesis}; as before, a rate becomes a rest energy only through the added
identification $mc^2=\hbar\kappa$.

\begin{table}[t]
\centering
\footnotesize
\begin{tabular}{@{}>{\raggedright\arraybackslash}p{3.6cm}>{\raggedright\arraybackslash}p{7.4cm}>{\raggedright\arraybackslash}p{3.4cm}@{}}
\toprule
statement & what the arena gives & verdict\\
\midrule
the algebra is duplicated & $\Cl(6,6)=\Cl(3,3)\hat\otimes\Cl(3,3)$, Klein factor, Fock space (\cref{d5:sec:algebra}) & forced; inherited\\
the energy is duplicated & a beam-splitter $S^3$ of covariant embeddings, one $\eta$-respecting circle; $z\mapsto(z,z)$ doubles the energy, $(z,z)/\sqrt2$ divides it; cloning of unknown quantum states excluded (\cref{d5:thm:dup}) & chosen; cloning obstructed (inherited)\\
the Lorentz action on the copies & one of five at $n=2$; it decides the spin of the duplicate and what an inter-manifold field reaches (\cref{d5:tab:actions,d5:prop:spin,d7:prop:exist}) & chosen\\
the duplicate adds structure & complex structure, copy charge, $\so{2,\C}$, $S^4$ vacuum, balance surface (\cref{d5:prop:cx,d5:prop:so6C,d5:prop:vacuum,d6:thm:balance}) & fixed (partly restated)\\
a pair from one null field & impossible covariantly; two null fields or a third body (\cref{d6:prop:covariant,d6:thm:pair}) & obstructed; inherited kinematics\\
what performs the duplication & not the closure (\cref{d6:cor:nopair}); the copy rotation: a gauge turn at equal rates, Rabi mixing at unequal ones (\cref{d6:sec:gauge}); a copy rotor (\cref{d6:rem:copyrotor}) & obstructed in the closure; mixing inherited; rotor constructed\\
the mass of one copy & not invariant under the copy rotation (\cref{d6:prop:gauge}) & fixed; a distinguishing structure open\\
a $b$-field between manifolds & $\so{n,\C}$, compact part $\so n$; a spectator manifold decouples (\cref{d7:prop:exist}) & fixed\\
it is a deal breaker & non-removable at a fixed rate matrix where it couples unequal rates (on an open interval $e^{-\tau b}$ removes it at the price of a time-dependent rate matrix): seesaw, confinement, resonance; a time-odd invariant at $n=3$ (\cref{d7:thm:schr,d7:thm:three}); removable in the common-rate closure (\cref{d7:rem:verdict}) & fixed\\
it generates a rate & rates are constant in the fixed-rate model, not by gauge invariance (\cref{d7:thm:schr}); closed worldlines need a holonomy statement & obstructed in the fixed-rate model\\
the rates differ & not in the common-rate closure; a non-central rotor in $\uu n$ admits them (\cref{d7:prop:rotor}) & constructed; origin open\\
isospin one-half between manifolds & copy isospin integral at $n=2,3$ under a common action; doublets from $n=4$ (\cref{d7:prop:isospin,d7:prop:so4unequal}) & obstructed at $n=2,3$\\
\bottomrule
\end{tabular}
\caption{Verdicts on duplication and on the $b$-field between manifolds (\cref{sec:dupsynthesis}).}
\label{tab:dupsynthesis}
\end{table}

\paragraph{How a unitary energy results in a duplication.} The algebra step is forced and the energy step is chosen. Doubling the vector space
forces the graded tensor product, whose spinor module is the Fock space of six fermionic modes. The undivided energy embeds isometrically in the
duplicate, and the Lorentz-covariant embeddings form the three-sphere of a lossless beam splitter, of which one circle respects $\eta$. Read as a
multiplication of the energy, the phrase fails: the well-defined classical map $z\mapsto(z,z)$ doubles the quadratic energy,
$E_{12}(z,z)=2E_6(z)$, while $z\mapsto(az,bz)$ with $|a|^2+|b|^2=1$, for instance $(z,z)/\sqrt2$, preserves it and divides it between the copies;
implementing an energy-increasing map physically requires an energy source. The no-cloning theorem, which concerns unknown quantum states, is a
separate statement about copying a quantum. The algebra step creates no pair: in $\Cl(6,6)$ the graded masses add in quadrature, and the pair threshold $2m$ appears
only for the ungraded operator (\cref{d6:prop:dirac}). What the duplicate can do that one copy cannot is precise: massive constituents of the
closure can turn at an exactly constant rate on a codimension-two balance surface (\cref{d6:thm:balance}), with no $b$-field.

\paragraph{The $b$-field between manifolds.} The remark that the $b$-field is a deal breaker between manifolds is true for transmission and false
for creation. A compact binding-compatible field coupling manifolds of unequal rates cannot be removed along a worldline while the rate matrix is kept fixed. On
an open interval the gauge $e^{-\tau b}$ removes it at the price of a time-dependent rate matrix, with $\tfrac{d}{d\tau}K'=[b,K]$ at $\tau=0$, and on
a closed worldline a holonomy statement is needed. At a fixed rate matrix it shifts and spreads
the lines, confines a frame without rate of its own by a seesaw, and from three manifolds on carries a time-odd invariant, the arena's analogue
of Jarlskog's invariant (\cref{d7:thm:schr,d7:thm:three}). But in the fixed-rate model the rates are the spectrum of the self-adjoint part of the shift, which transforms
as a tensor, so no field creates or changes a rate; that the rates are constant follows from the model, not from gauge invariance; with equal rates every binding-compatible field is removable, $\so4_\Delta$ included, and in the
common-rate closure the deal breaker has no effect (\cref{d7:rem:verdict}). The part of the copy field that survives every change of frame is
proportional to the rate differences. Read with $mc^2=\hbar\kappa$, the order of the hypothesis is reversed: the differentiation of the rates
breaks the copy symmetry, and the $b$-field acts only because it is broken. Formally this has the structure of the Englert--Brout--Higgs--Kibble
mechanism with the rate matrix as an adjoint vacuum value; that reading is interpretive, since here $b$ has no kinetic term and the rates are
inputs. Their origin is open: a non-central rotor generator, admitted but not selected, or a quantum scale by transmutation, which is
conjectural.
